\documentclass[a4paper,12pt]{article}

% Packages
\usepackage[utf8]{inputenc}
\usepackage[T1]{fontenc}
\usepackage{mathptmx} % Times New Roman font
\usepackage{listings} % For code display
\usepackage{graphicx} % For images
\usepackage{fancyhdr} % For headers and footers
\usepackage{geometry} % For page layout

% Page geometry
\geometry{
    a4paper,
    left=1in,
    right=1in,
    top=1in,
    bottom=1in
}

% Line spacing
\usepackage{setspace}
\onehalfspacing

% Header and Footer setup
\pagestyle{fancy}
\fancyhf{}
\fancyhead[L]{Cryptography and Network Security(3161606)}
\fancyhead[R]{Enrollment No.}
\fancyfoot[C]{\thepage}

% Code listing style

\begin{document}

% Title
\begin{center}
    \Large\textbf{Practical No.8}
\end{center}

% Aim
\section*{Aim}
\large\textbf{To implement and demonstrate the RSA algorithm for encryption and decryption using modular arithmetic.}
\normalsize

% Theory
\section*{Theory}
\textbf{RSA Algorithm} is a public-key cryptosystem that uses modular arithmetic for secure data transmission. The algorithm involves:
\begin{itemize}
    \item \textbf{Key Generation:} Select two large prime numbers p and q, compute n = p × q and φ(n) = (p-1) × (q-1). Choose public key e such that 1 < e < φ(n) and gcd(e, φ(n)) = 1. Calculate private key d such that d × e ≡ 1 (mod φ(n)).
    \item \textbf{Encryption:} Ciphertext C = M$^e$ mod n, where M is the plaintext.
    \item \textbf{Decryption:} Plaintext M = C$^d$ mod n, where C is the ciphertext.
\end{itemize}

% Code
\section*{Code}
\begin{verbatim}
/**
 * RSA Algorithm Implementation - Practical 8
 * Test case: p=17, q=11, e=7, M=88, Expected: d=23, C=11
 */

// ============ MODULAR ARITHMETIC FUNCTIONS ============

function modExp(base: bigint, exp: bigint, mod: bigint): bigint {
  let result = 1n
  base = base % mod
  while (exp > 0n) {
    if (exp % 2n === 1n) result = (result * base) % mod
    exp = exp / 2n
    base = (base * base) % mod
  }
  return result
}

function gcd(a: bigint, b: bigint): bigint {
  while (b !== 0n) {
    const temp = b
    b = a % b
    a = temp
  }
  return a
}

function extendedGcd(
  a: bigint,
  b: bigint
): [bigint, bigint, bigint] {
  if (a === 0n) return [b, 0n, 1n]
  const [g, x1, y1] = extendedGcd(b % a, a)
  const x = y1 - (b / a) * x1
  const y = x1
  return [g, x, y]
}

function modInverse(a: bigint, m: bigint): bigint | null {
  const [g, x] = extendedGcd(a, m)
  if (g !== 1n) return null
  return ((x % m) + m) % m
}

// ============ KEY GENERATION ============

function generateRSAKeyPair(p: bigint, q: bigint, e: bigint) {
  const n = p * q
  const phi = (p - 1n) * (q - 1n)
  const d = modInverse(e, phi)

  return {
    publicKey: { n, e },
    privateKey: { n, d },
    phi
  }
}

// ============ ENCRYPTION & DECRYPTION ============

function rsaEncrypt(M: bigint, publicKey: { n: bigint; e: bigint }): bigint {
  return modExp(M, publicKey.e, publicKey.n)
}

function rsaDecrypt(C: bigint, privateKey: { n: bigint; d: bigint }): bigint {
  return modExp(C, privateKey.d, privateKey.n)
}

// ============ TEST CASE ============

console.log('='.repeat(70))
console.log('RSA ALGORITHM - PRACTICAL 8')
console.log('='.repeat(70))

// Test case parameters
const p = 17n
const q = 11n
const e = 7n
const M = 88n
const expectedD = 23n
const expectedC = 11n

console.log('\n--- Test Case Parameters ---')
console.log('Prime p = ' + p)
console.log('Prime q = ' + q)
console.log('Encryption key e = ' + e)
console.log('Plaintext M = ' + M)

// Key generation
console.log('\n--- Key Generation ---')
const keyPair = generateRSAKeyPair(p, q, e)
console.log('n = p × q = ' + keyPair.publicKey.n)
console.log('φ(n) = (p-1) × (q-1) = ' + keyPair.phi)
console.log('e = ' + keyPair.publicKey.e)
console.log('d = ' + keyPair.privateKey.d)
console.log('Verification: d × e mod φ(n) = ' + ((keyPair.privateKey.d * keyPair.publicKey.e) % keyPair.phi))

// Verify decryption key
console.log('\n--- Verification ---')
if (keyPair.privateKey.d === expectedD) {
  console.log('✓ Decryption key verified: d = ' + keyPair.privateKey.d)
} else {
  console.log(
    '✗ Decryption key mismatch: expected ' +
      expectedD +
      ', got ' +
      keyPair.privateKey.d
  )
}

// Encryption
console.log('\n--- Encryption ---')
const C = rsaEncrypt(M, keyPair.publicKey)
console.log('C = M^e mod n')
console.log('C = ' + M + '^' + keyPair.publicKey.e + ' mod ' + keyPair.publicKey.n)
console.log('C = ' + C)

// Verify ciphertext
if (C === expectedC) {
  console.log('✓ Ciphertext verified: C = ' + C)
} else {
  console.log('✗ Ciphertext mismatch: expected ' + expectedC + ', got ' + C)
}

// Decryption
console.log('\n--- Decryption ---')
const decryptedM = rsaDecrypt(C, keyPair.privateKey)
console.log('M = C^d mod n')
console.log('M = ' + C + '^' + keyPair.privateKey.d + ' mod ' + keyPair.privateKey.n)
console.log('M = ' + decryptedM)

// Verify decryption
if (decryptedM === M) {
  console.log('✓ Decryption verified: M = ' + decryptedM)
} else {
  console.log('✗ Decryption failed: expected ' + M + ', got ' + decryptedM)
}

// Summary
console.log('\n' + '='.repeat(70))
console.log('SUMMARY')
console.log('='.repeat(70))
console.log('Test case: ' + (C === expectedC && decryptedM === M ? 'PASSED ✓' : 'FAILED ✗'))
console.log('All verifications: ' + (keyPair.privateKey.d === expectedD && C === expectedC && decryptedM === M ? 'PASSED ✓' : 'FAILED ✗'))
console.log('='.repeat(70))
\end{verbatim}

% Output
\section*{Output}
\begin{figure}[h]
    \centering
    \includegraphics[width=\textwidth]{practical8-output.png}
    \caption{RSA Algorithm Implementation Output}
\end{figure}

\end{document}
